
%%%%%%%%%%%%%%%%%%%%%%%%%%%%%%%%%%%%%%%%%%%%%%%%%%%%%%%%%%%%%%%%%%%%%%%%%%%%%%%%%%%%%%%%%%%%%
\begin{frame}[allowframebreaks,c]{} \label{MLFF}

\begin{center}
\Huge \textbf{MLFF}
\end{center}

\end{frame}
%%%%%%%%%%%%%%%%%%%%%%%%%%%%%%%%%%%%%%%%%%%%%%%%%%%%%%%%%%%%%%%%%%%%%%%%%%%%%%%%%%%%%%%%%%%%%

%%%%%%%%%%%%%%%%%%%%%%%%%%%%%%%%%%%%%%%%%%%%%%%%%%%%%%%%%%%%%%%%%%%%%%%%%%%%%%%%%%%%%%%%%%%%%
\begin{frame}[allowframebreaks]{\texttt{MLFF\_FLAG}} \label{MLFF_FLAG}
\vspace*{-12pt}
\begin{columns}
\column{0.4\linewidth}
\begin{block}{Type}
Integer
\end{block}

\begin{block}{Default}
0
\end{block}

\column{0.4\linewidth}
\begin{block}{Unit}
No unit
\end{block}

\begin{block}{Example}
\texttt{MLFF\_FLAG}: 1
\end{block}
\end{columns}

\begin{block}{Description}
Flag to turn on MLFF in SPARC. There are three options currently available. (1) MLFF\_FLAG: 1 to perform on-the-fly MD with no existing model, (2) MLFF\_FLAG: 21 to perform only predictions from an existing model, and (3) MLFF\_FLAG: 22 to perform on-the-fly MD building on top of an existing model
\end{block}

\end{frame}
%%%%%%%%%%%%%%%%%%%%%%%%%%%%%%%%%%%%%%%%%%%%%%%%%%%%%%%%%%%%%%%%%%%%%%%%%%%%%%%%%%%%%%%%%%%%%





%%%%%%%%%%%%%%%%%%%%%%%%%%%%%%%%%%%%%%%%%%%%%%%%%%%%%%%%%%%%%%%%%%%%%%%%%%%%%%%%%%%%%%%%%%%%%
\begin{frame}[allowframebreaks]{\texttt{{MLFF\_RCUT\_SOAP}}} \label{MLFF_RCUT_SOAP}
\vspace*{-12pt}
\begin{columns}
\column{0.4\linewidth}
\begin{block}{Type}
Double
\end{block}

\begin{block}{Default}
10.0
\end{block}

\column{0.4\linewidth}
\begin{block}{Unit}
Bohr
\end{block}

\begin{block}{Example}
\texttt{MLFF\_RCUT\_SOAP}: 8.0
\end{block}
\end{columns}

\begin{block}{Description}
 The cutoff radius used to calculate the SOAP descriptor. 
\end{block}

%\begin{block}{Remark}
%
%\end{block}
\end{frame}
%%%%%%%%%%%%%%%%%%%%%%%%%%%%%%%%%%%%%%%%%%%%%%%%%%%%%%%%%%%%%%%%%%%%%%%%%%%%%%%%%%%%%%%%%%%%%

%%%%%%%%%%%%%%%%%%%%%%%%%%%%%%%%%%%%%%%%%%%%%%%%%%%%%%%%%%%%%%%%%%%%%%%%%%%%%%%%%%%%%%%%%%%%%
\begin{frame}[allowframebreaks]{\texttt{{MLFF\_SIGMA\_ATOM\_SOAP}}} \label{MLFF_SIGMA_ATOM_SOAP}
\vspace*{-12pt}
\begin{columns}
\column{0.4\linewidth}
\begin{block}{Type}
Double
\end{block}

\begin{block}{Default}
1.0
\end{block}

\column{0.4\linewidth}
\begin{block}{Unit}
Bohr
\end{block}

\begin{block}{Example}
\texttt{MLFF\_SIGMA\_ATOM\_SOAP}: 1.0
\end{block}
\end{columns}

\begin{block}{Description}
 The width of Gaussians in SOAP descriptor. 
\end{block}

%\begin{block}{Remark}
%
%\end{block}
\end{frame}
%%%%%%%%%%%%%%%%%%%%%%%%%%%%%%%%%%%%%%%%%%%%%%%%%%%%%%%%%%%%%%%%%%%%%%%%%%%%%%%%%%%%%%%%%%%%%

%%%%%%%%%%%%%%%%%%%%%%%%%%%%%%%%%%%%%%%%%%%%%%%%%%%%%%%%%%%%%%%%%%%%%%%%%%%%%%%%%%%%%%%%%%%%%
\begin{frame}[allowframebreaks]{\texttt{{MLFF\_RADIAL\_BASIS}}} \label{MLFF_RADIAL_BASIS}
\vspace*{-12pt}
\begin{columns}
\column{0.4\linewidth}
\begin{block}{Type}
Integer
\end{block}

\begin{block}{Default}
8
\end{block}

\column{0.4\linewidth}
\begin{block}{Unit}
No unit
\end{block}

\begin{block}{Example}
\texttt{MLFF\_RADIAL\_BASIS}: 7
\end{block}
\end{columns}

\begin{block}{Description}
 The number of radial basis functions used in SOAP descriptors.
\end{block}

\begin{block}{Remark}
8-10 radial basis functions should be sufficient for most applications.
\end{block}
\end{frame}
%%%%%%%%%%%%%%%%%%%%%%%%%%%%%%%%%%%%%%%%%%%%%%%%%%%%%%%%%%%%%%%%%%%%%%%%%%%%%%%%%%%%%%%%%%%%%



%%%%%%%%%%%%%%%%%%%%%%%%%%%%%%%%%%%%%%%%%%%%%%%%%%%%%%%%%%%%%%%%%%%%%%%%%%%%%%%%%%%%%%%%%%%%%
\begin{frame}[allowframebreaks]{\texttt{{MLFF\_ANGULAR\_BASIS}}} \label{MLFF_ANGULAR_BASIS}
\vspace*{-12pt}
\begin{columns}
\column{0.4\linewidth}
\begin{block}{Type}
Integer
\end{block}

\begin{block}{Default}
6
\end{block}

\column{0.4\linewidth}
\begin{block}{Unit}
No unit
\end{block}

\begin{block}{Example}
\texttt{MLFF\_ANGULAR\_BASIS}: 7
\end{block}
\end{columns}

\begin{block}{Description}
 The maximum angular momentum quantum number (L) of spherical harmonics used in the SOAP descriptor.
\end{block}

\begin{block}{Remark}
6-8 for SOAP works well for most applications.
\end{block}
\end{frame}
%%%%%%%%%%%%%%%%%%%%%%%%%%%%%%%%%%%%%%%%%%%%%%%%%%%%%%%%%%%%%%%%%%%%%%%%%%%%%%%%%%%%%%%%%%%%%

%%%%%%%%%%%%%%%%%%%%%%%%%%%%%%%%%%%%%%%%%%%%%%%%%%%%%%%%%%%%%%%%%%%%%%%%%%%%%%%%%%%%%%%%%%%%%
\begin{frame}[allowframebreaks]{\texttt{{MLFF\_INITIAL\_STEPS\_TRAIN}}} \label{MLFF_INITIAL_STEPS_TRAIN}
\vspace*{-12pt}
\begin{columns}
\column{0.4\linewidth}
\begin{block}{Type}
Integer
\end{block}

\begin{block}{Default}
10
\end{block}

\column{0.4\linewidth}
\begin{block}{Unit}
No unit
\end{block}

\begin{block}{Example}
\texttt{MLFF\_INITIAL\_STEPS\_TRAIN}: 20
\end{block}
\end{columns}

\begin{block}{Description}
 The number of DFT steps performed in the beginning of on-the-fly MLFF simulation.
\end{block}

%\begin{block}{Remark}
%
%\end{block}
\end{frame}
%%%%%%%%%%%%%%%%%%%%%%%%%%%%%%%%%%%%%%%%%%%%%%%%%%%%%%%%%%%%%%%%%%%%%%%%%%%%%%%%%%%%%%%%%%%%%

%%%%%%%%%%%%%%%%%%%%%%%%%%%%%%%%%%%%%%%%%%%%%%%%%%%%%%%%%%%%%%%%%%%%%%%%%%%%%%%%%%%%%%%%%%%%%
\begin{frame}[allowframebreaks]{\texttt{{MLFF\_MODEL\_FOLDER}}} \label{MLFF_MODEL_FOLDER}
\vspace*{-12pt}
\begin{columns}
\column{0.4\linewidth}
\begin{block}{Type}
String
\end{block}

\begin{block}{Default}
No default
\end{block}

\column{0.4\linewidth}
\begin{block}{Unit}
No unit
\end{block}

\begin{block}{Example}
\texttt{MLFF\_MODEL\_FOLDER}: ././
\end{block}
\end{columns}

\begin{block}{Description}
 The location of folder where the trained model file is located.
\end{block}

%\begin{block}{Remark}
%
%\end{block}
\end{frame}
%%%%%%%%%%%%%%%%%%%%%%%%%%%%%%%%%%%%%%%%%%%%%%%%%%%%%%%%%%%%%%%%%%%%%%%%%%%%%%%%%%%%%%%%%%%%%



%%%%%%%%%%%%%%%%%%%%%%%%%%%%%%%%%%%%%%%%%%%%%%%%%%%%%%%%%%%%%%%%%%%%%%%%%%%%%%%%%%%%%%%%%%%%%
\begin{frame}[allowframebreaks]{\texttt{{MLFF\_IF\_ATOM\_DATA\_AVAILABLE}}} \label{MLFF_IF_ATOM_DATA_AVAILABLE}
\vspace*{-12pt}
\begin{columns}
\column{0.4\linewidth}
\begin{block}{Type}
Integer
\end{block}

\begin{block}{Default}
0
\end{block}

\column{0.4\linewidth}
\begin{block}{Unit}
No unit
\end{block}

\begin{block}{Example}
\texttt{MLFF\_IF\_ATOM\_DATA\_AVAILABLE}: 1
\end{block}
\end{columns}

\begin{block}{Description}
 Flag to indicate if an pretrained MLFF model is available. 
\end{block}

%\begin{block}{Remark}
%
%\end{block}
\end{frame}
%%%%%%%%%%%%%%%%%%%%%%%%%%%%%%%%%%%%%%%%%%%%%%%%%%%%%%%%%%%%%%%%%%%%%%%%%%%%%%%%%%%%%%%%%%%%%

%%%%%%%%%%%%%%%%%%%%%%%%%%%%%%%%%%%%%%%%%%%%%%%%%%%%%%%%%%%%%%%%%%%%%%%%%%%%%%%%%%%%%%%%%%%%%
\begin{frame}[allowframebreaks]{\texttt{{MLFF\_REGUL\_MIN}}} \label{MLFF_REGUL_MIN}
\vspace*{-12pt}
\begin{columns}
\column{0.4\linewidth}
\begin{block}{Type}
Double
\end{block}

\begin{block}{Default}
1e-14
\end{block}

\column{0.4\linewidth}
\begin{block}{Unit}
No unit
\end{block}

\begin{block}{Example}
\texttt{MLFF\_REGUL\_MIN}: 1e-10
\end{block}
\end{columns}

\begin{block}{Description}
 The mimimum value of inverse of condition number of $K^TK + \lambda I$ matrix which appear in MLFF calculation.
\end{block}

\begin{block}{Remark}
The $K^TK + \lambda I$ matrix needs to be inverted during the training of MLFF. The matrix $K^TK$ is in general ill-conditioned so regularization is used to improve the conditioning. Any number in the range of 1E-10-1E-14 should work. A larger value should result in lesser DFT calls, but also lesser accuracy of MLFF model.
\end{block}
\end{frame}
%%%%%%%%%%%%%%%%%%%%%%%%%%%%%%%%%%%%%%%%%%%%%%%%%%%%%%%%%%%%%%%%%%%%%%%%%%%%%%%%%%%%%%%%%%%%%

%%%%%%%%%%%%%%%%%%%%%%%%%%%%%%%%%%%%%%%%%%%%%%%%%%%%%%%%%%%%%%%%%%%%%%%%%%%%%%%%%%%%%%%%%%%%%
\begin{frame}[allowframebreaks]{\texttt{{MLFF\_MAX\_STR\_STORE}}} \label{MLFF_MAX_STR_STORE}
\vspace*{-12pt}
\begin{columns}
\column{0.4\linewidth}
\begin{block}{Type}
Integer
\end{block}

\begin{block}{Default}
500
\end{block}

\column{0.4\linewidth}
\begin{block}{Unit}
No unit
\end{block}

\begin{block}{Example}
\texttt{MLFF\_MAX\_STR\_STORE}: 1000
\end{block}
\end{columns}

\begin{block}{Description}
 Maximum number of DFT calls that can be made during the on-the-fly simulation. 
\end{block}

\begin{block}{Remark}
A larger number requires larger memory allocation.
\end{block}
\end{frame}
%%%%%%%%%%%%%%%%%%%%%%%%%%%%%%%%%%%%%%%%%%%%%%%%%%%%%%%%%%%%%%%%%%%%%%%%%%%%%%%%%%%%%%%%%%%%%

%%%%%%%%%%%%%%%%%%%%%%%%%%%%%%%%%%%%%%%%%%%%%%%%%%%%%%%%%%%%%%%%%%%%%%%%%%%%%%%%%%%%%%%%%%%%%
\begin{frame}[allowframebreaks]{\texttt{{MLFF\_MAX\_CONFIG\_STORE}}} \label{MLFF_MAX_CONFIG_STORE}
\vspace*{-12pt}
\begin{columns}
\column{0.4\linewidth}
\begin{block}{Type}
Integer
\end{block}

\begin{block}{Default}
5000
\end{block}

\column{0.4\linewidth}
\begin{block}{Unit}
No unit
\end{block}

\begin{block}{Example}
\texttt{MLFF\_MAX\_CONFIG\_STORE}: 1000
\end{block}
\end{columns}

\begin{block}{Description}
 Maximum number of atomic descriptors per element type that will be used in training.
\end{block}

\begin{block}{Remark}
A larger number requires larger memory allocation.
\end{block}
\end{frame}
%%%%%%%%%%%%%%%%%%%%%%%%%%%%%%%%%%%%%%%%%%%%%%%%%%%%%%%%%%%%%%%%%%%%%%%%%%%%%%%%%%%%%%%%%%%%%


%%%%%%%%%%%%%%%%%%%%%%%%%%%%%%%%%%%%%%%%%%%%%%%%%%%%%%%%%%%%%%%%%%%%%%%%%%%%%%%%%%%%%%%%%%%%%
\begin{frame}[allowframebreaks]{\texttt{{MLFF\_FACTOR\_MULTIPLY\_SIGMATOL}}} \label{MLFF_FACTOR_MULTIPLY_SIGMATOL}
\vspace*{-12pt}
\begin{columns}
\column{0.4\linewidth}
\begin{block}{Type}
Double
\end{block}

\begin{block}{Default}
1.01
\end{block}

\column{0.4\linewidth}
\begin{block}{Unit}
No unit
\end{block}

\begin{block}{Example}
\texttt{MLFF\_FACTOR\_MULTIPLY\_SIGMATOL}: 1.001
\end{block}
\end{columns}

\begin{block}{Description}
 A prefactor multiplied with maximum of Bayesian errors while updating the $\sigma_{\rm tol}$. It has to be greater than 1.
\end{block}

\begin{block}{Remark}
A larger value will result lesser DFT calls hence might reduce the accuracy. Use any number in (1.0001-1.01) range.
\end{block}
\end{frame}
%%%%%%%%%%%%%%%%%%%%%%%%%%%%%%%%%%%%%%%%%%%%%%%%%%%%%%%%%%%%%%%%%%%%%%%%%%%%%%%%%%%%%%%%%%%%%

%%%%%%%%%%%%%%%%%%%%%%%%%%%%%%%%%%%%%%%%%%%%%%%%%%%%%%%%%%%%%%%%%%%%%%%%%%%%%%%%%%%%%%%%%%%%%
\begin{frame}[allowframebreaks]{\texttt{{MLFF\_IF\_SPARSIFY\_BEFORE\_TRAIN}}} \label{MLFF_IF_SPARSIFY_BEFORE_TRAIN}
\vspace*{-12pt}
\begin{columns}
\column{0.4\linewidth}
\begin{block}{Type}
Integer
\end{block}

\begin{block}{Default}
1
\end{block}

\column{0.4\linewidth}
\begin{block}{Unit}
No unit
\end{block}

\begin{block}{Example}
\texttt{MLFF\_IF\_SPARSIFY\_BEFORE\_TRAIN}: 1
\end{block}
\end{columns}

\begin{block}{Description}
 Flag to turn on the sparsification before each of the training steps.
\end{block}

%\begin{block}{Remark}
%
%\end{block}
\end{frame}
%%%%%%%%%%%%%%%%%%%%%%%%%%%%%%%%%%%%%%%%%%%%%%%%%%%%%%%%%%%%%%%%%%%%%%%%%%%%%%%%%%%%%%%%%%%%%

%%%%%%%%%%%%%%%%%%%%%%%%%%%%%%%%%%%%%%%%%%%%%%%%%%%%%%%%%%%%%%%%%%%%%%%%%%%%%%%%%%%%%%%%%%%%%
\begin{frame}[allowframebreaks]{\texttt{{MLFF\_EXPONENT\_SOAP}}} \label{MLFF_EXPONENT_SOAP}
\vspace*{-12pt}
\begin{columns}
\column{0.4\linewidth}
\begin{block}{Type}
Double
\end{block}

\begin{block}{Default}
4.0
\end{block}

\column{0.4\linewidth}
\begin{block}{Unit}
No unit
\end{block}

\begin{block}{Example}
\texttt{MLFF\_EXPONENT\_SOAP}: 4.0
\end{block}
\end{columns}

\begin{block}{Description}
 The exponent in the polynomial kernel used in Bayesian linear regression.
\end{block}

\begin{block}{Remark}
In principle, any number greater than 1 could be used. The code has been tested most for the exponent of 4.
\end{block}
\end{frame}
%%%%%%%%%%%%%%%%%%%%%%%%%%%%%%%%%%%%%%%%%%%%%%%%%%%%%%%%%%%%%%%%%%%%%%%%%%%%%%%%%%%%%%%%%%%%%

%%%%%%%%%%%%%%%%%%%%%%%%%%%%%%%%%%%%%%%%%%%%%%%%%%%%%%%%%%%%%%%%%%%%%%%%%%%%%%%%%%%%%%%%%%%%%
\begin{frame}[allowframebreaks]{\texttt{{MLFF\_TOL\_FORCE}}} \label{MLFF_TOL_FORCE}
\vspace*{-12pt}
\begin{columns}
\column{0.4\linewidth}
\begin{block}{Type}
Double
\end{block}

\begin{block}{Default}
5e-10
\end{block}

\column{0.4\linewidth}
\begin{block}{Unit}
No unit
\end{block}

\begin{block}{Example}
\texttt{MLFF\_TOL\_FORCE}: 1e-10
\end{block}
\end{columns}

\begin{block}{Description}
 The initial value of $\sigma_{\rm tol}$.
\end{block}

%\begin{block}{Remark}
%
%\end{block}
\end{frame}
%%%%%%%%%%%%%%%%%%%%%%%%%%%%%%%%%%%%%%%%%%%%%%%%%%%%%%%%%%%%%%%%%%%%%%%%%%%%%%%%%%%%%%%%%%%%%

%%%%%%%%%%%%%%%%%%%%%%%%%%%%%%%%%%%%%%%%%%%%%%%%%%%%%%%%%%%%%%%%%%%%%%%%%%%%%%%%%%%%%%%%%%%%%
\begin{frame}[allowframebreaks]{\texttt{{MLFF\_SCALE\_FORCE}}} \label{MLFF_SCALE_FORCE}
\vspace*{-12pt}
\begin{columns}
\column{0.4\linewidth}
\begin{block}{Type}
Double
\end{block}

\begin{block}{Default}
1.0
\end{block}

\column{0.4\linewidth}
\begin{block}{Unit}
No unit
\end{block}

\begin{block}{Example}
\texttt{MLFF\_SCALE\_FORCE}: 10
\end{block}
\end{columns}

\begin{block}{Description}
 The scaling of force terms in the loss function during the training.
\end{block}

\begin{block}{Remark}
A larger number can result in better accuracy of MLFF forces.
\end{block}
\end{frame}
%%%%%%%%%%%%%%%%%%%%%%%%%%%%%%%%%%%%%%%%%%%%%%%%%%%%%%%%%%%%%%%%%%%%%%%%%%%%%%%%%%%%%%%%%%%%%

%%%%%%%%%%%%%%%%%%%%%%%%%%%%%%%%%%%%%%%%%%%%%%%%%%%%%%%%%%%%%%%%%%%%%%%%%%%%%%%%%%%%%%%%%%%%%
\begin{frame}[allowframebreaks]{\texttt{{MLFF\_SPLINE\_NGRID}}} \label{MLFF_SPLINE_NGRID}
\vspace*{-12pt}
\begin{columns}
\column{0.4\linewidth}
\begin{block}{Type}
Integer
\end{block}

\begin{block}{Default}
100
\end{block}

\column{0.4\linewidth}
\begin{block}{Unit}
No unit
\end{block}

\begin{block}{Example}
\texttt{MLFF\_SPLINE\_NGRID}: 100
\end{block}
\end{columns}

\begin{block}{Description}
 Number of grid points to use for spline interpolation of hnl.
\end{block}

\begin{block}{Remark}
A larger number can result in better accuracy of MLFF forces.
\end{block}
\end{frame}
%%%%%%%%%%%%%%%%%%%%%%%%%%%%%%%%%%%%%%%%%%%%%%%%%%%%%%%%%%%%%%%%%%%%%%%%%%%%%%%%%%%%%%%%%%%%%

%%%%%%%%%%%%%%%%%%%%%%%%%%%%%%%%%%%%%%%%%%%%%%%%%%%%%%%%%%%%%%%%%%%%%%%%%%%%%%%%%%%%%%%%%%%%%
\begin{frame}[allowframebreaks]{\texttt{{MLFF\_SCALE\_STRESS}}} \label{MLFF_SCALE_STRESS}
\vspace*{-12pt}
\begin{columns}
\column{0.4\linewidth}
\begin{block}{Type}
Double
\end{block}

\begin{block}{Default}
1.0 1.0 1.0 1.0 1.0 1.0
\end{block}

\column{0.4\linewidth}
\begin{block}{Unit}
No unit
\end{block}

\begin{block}{Example}
\texttt{MLFF\_SCALE\_STRESS}: 10 1 1 10 1 10
\end{block}
\end{columns}

\begin{block}{Description}
 The scaling of stress tensor terms in the loss function during the training.
\end{block}

%\begin{block}{Remark}
%
%\end{block}
\end{frame}
%%%%%%%%%%%%%%%%%%%%%%%%%%%%%%%%%%%%%%%%%%%%%%%%%%%%%%%%%%%%%%%%%%%%%%%%%%%%%%%%%%%%%%%%%%%%%

%%%%%%%%%%%%%%%%%%%%%%%%%%%%%%%%%%%%%%%%%%%%%%%%%%%%%%%%%%%%%%%%%%%%%%%%%%%%%%%%%%%%%%%%%%%%%
\begin{frame}[allowframebreaks]{\texttt{{MLFF\_PRINT\_FLAG}}} \label{MLFF_PRINT_FLAG}
\vspace*{-12pt}
\begin{columns}
\column{0.4\linewidth}
\begin{block}{Type}
Integer
\end{block}

\begin{block}{Default}
1
\end{block}

\column{0.4\linewidth}
\begin{block}{Unit}
No unit
\end{block}

\begin{block}{Example}
\texttt{MLFF\_PRINT\_FLAG}: 1
\end{block}
\end{columns}

\begin{block}{Description}
Flag to turn on the printing of MLFF related log in a file named 'mlff.log' during runtime.
\end{block}

\end{frame}
%%%%%%%%%%%%%%%%%%%%%%%%%%%%%%%%%%%%%%%%%%%%%%%%%%%%%%%%%%%%%%%%%%%%%%%%%%%%%%%%%%%%%%%%%%%%%

%%%%%%%%%%%%%%%%%%%%%%%%%%%%%%%%%%%%%%%%%%%%%%%%%%%%%%%%%%%%%%%%%%%%%%%%%%%%%%%%%%%%%%%%%%%%%
\begin{frame}[allowframebreaks]{\texttt{{MLFF\_INTERNAL\_ENERGY\_FLAG}}} \label{MLFF_INTERNAL_ENERGY_FLAG}
\vspace*{-12pt}
\begin{columns}
\column{0.4\linewidth}
\begin{block}{Type}
Integer
\end{block}

\begin{block}{Default}
0
\end{block}

\column{0.4\linewidth}
\begin{block}{Unit}
No unit
\end{block}

\begin{block}{Example}
\texttt{MLFF\_INTERNAL\_ENERGY\_FLAG}: 1
\end{block}
\end{columns}

\begin{block}{Description}
Flag to turn on the training of internal energy model based on linear regression.
\end{block}

\end{frame}
%%%%%%%%%%%%%%%%%%%%%%%%%%%%%%%%%%%%%%%%%%%%%%%%%%%%%%%%%%%%%%%%%%%%%%%%%%%%%%%%%%%%%%%%%%%%%

%%%%%%%%%%%%%%%%%%%%%%%%%%%%%%%%%%%%%%%%%%%%%%%%%%%%%%%%%%%%%%%%%%%%%%%%%%%%%%%%%%%%%%%%%%%%%
\begin{frame}[allowframebreaks]{\texttt{{MLFF\_PRESSURE\_TRAIN\_FLAG}}} \label{MLFF_PRESSURE_TRAIN_FLAG}
\vspace*{-12pt}
\begin{columns}
\column{0.4\linewidth}
\begin{block}{Type}
Integer
\end{block}

\begin{block}{Default}
0
\end{block}

\column{0.4\linewidth}
\begin{block}{Unit}
No unit
\end{block}

\begin{block}{Example}
\texttt{MLFF\_PRESSURE\_TRAIN\_FLAG}: 1
\end{block}
\end{columns}

\begin{block}{Description}
Flag to turn on the training with only pressure and not all of the stress tensor components.
\end{block}

\end{frame}
%%%%%%%%%%%%%%%%%%%%%%%%%%%%%%%%%%%%%%%%%%%%%%%%%%%%%%%%%%%%%%%%%%%%%%%%%%%%%%%%%%%%%%%%%%%%%

\begin{frame}[allowframebreaks]{\texttt{{MLFF\_DFT\_FQ}}} \label{MLFF_DFT_FQ}
\vspace*{-12pt}
\begin{columns}
\column{0.4\linewidth}
\begin{block}{Type}
Integer
\end{block}

\begin{block}{Default}
100000000
\end{block}

\column{0.4\linewidth}
\begin{block}{Unit}
No unit
\end{block}

\begin{block}{Example}
\texttt{MLFF\_DFT\_FQ}: 100
\end{block}
\end{columns}

\begin{block}{Description}
Input variable to set the frequency at which DFT must be performed during on-the-fly MD.
\end{block}

\end{frame}
%%%%%%%%%%%%%%%%%%%%%%%%%%%%%%%%%%%%%%%%%%%%%%%%%%%%%%%%%%%%%%%%%%%%%%%%%%%%%%%%%%%%%%%%%%%%%

%%%%%%%%%%%%%%%%%%%%%%%%%%%%%%%%%%%%%%%%%%%%%%%%%%%%%%%%%%%%%%%%%%%%%%%%%%%%%%%%%%%%%%%%%%%%%
\begin{frame}[allowframebreaks]{\texttt{{MLFF\_SWITCH\_STEP}}} \label{MLFF_SWITCH_STEP}
\vspace*{-12pt}
\begin{columns}
\column{0.4\linewidth}
\begin{block}{Type}
Integer
\end{block}

\begin{block}{Default}
100
\end{block}

\column{0.4\linewidth}
\begin{block}{Unit}
No unit
\end{block}

\begin{block}{Example}
\texttt{MLFF\_SWITCH\_STEP}: 500
\end{block}
\end{columns}

\begin{block}{Description}
Input variable to set the switch over step after which the initial coarse MLFF model is deleted and the production on-the-fly MD starts.
\end{block}

\begin{block}{Remark}
This is relevant only for starting a on-the-fly MLFF MD from scratch.
\end{block}

\end{frame}
%%%%%%%%%%%%%%%%%%%%%%%%%%%%%%%%%%%%%%%%%%%%%%%%%%%%%%%%%%%%%%%%%%%%%%%%%%%%%%%%%%%%%%%%%%%%%

%%%%%%%%%%%%%%%%%%%%%%%%%%%%%%%%%%%%%%%%%%%%%%%%%%%%%%%%%%%%%%%%%%%%%%%%%%%%%%%%%%%%%%%%%%%%%
\begin{frame}[allowframebreaks]{\texttt{{MLFF\_FACTOR\_MULTIPLY\_SIGMATOL\_COARSE}}} \label{MLFF_FACTOR_MULTIPLY_SIGMATOL_COARSE}
\vspace*{-12pt}
\begin{columns}
\column{0.4\linewidth}
\begin{block}{Type}
Double
\end{block}

\column{0.4\linewidth}
\begin{block}{Unit}
No unit
\end{block}

\end{columns}

\begin{block}{Default}
$1.0 + 5\times(\texttt{{MLFF\_FACTOR\_MULTIPLY\_SIGMATOL}} - 1.0)$
\end{block}

\begin{block}{Example}
\texttt{MLFF\_FACTOR\_MULTIPLY\_SIGMATOL\_COARSE}: 1.05
\end{block}

\begin{block}{Description}
Prefactor multiplied with maximum Bayesian error while updating $\sigma_{\text{tol}}$ for the initial coarse model. Ideally, it is greater than \texttt{MLFF\_FACTOR\_MULTIPLY\_SIGMATOL} to skip more DFT steps. This is relevant only for starting a on-the-fly MLFF MD from scratch.
\end{block}

\end{frame}
%%%%%%%%%%%%%%%%%%%%%%%%%%%%%%%%%%%%%%%%%%%%%%%%%%%%%%%%%%%%%%%%%%%%%%%%%%%%%%%%%%%%%%%%%%%%%

%%%%%%%%%%%%%%%%%%%%%%%%%%%%%%%%%%%%%%%%%%%%%%%%%%%%%%%%%%%%%%%%%%%%%%%%%%%%%%%%%%%%%%%%%%%%%
\begin{frame}[allowframebreaks]{\texttt{{MLFF\_DELTA\_FLAG}}} \label{MLFF_DELTA_FLAG}
\vspace*{-12pt}
\begin{columns}
\column{0.4\linewidth}
\begin{block}{Type}
Integer
\end{block}

\begin{block}{Default}
0
\end{block}

\column{0.4\linewidth}
\begin{block}{Unit}
No unit
\end{block}

\begin{block}{Example}
\texttt{MLFF\_DELTA\_FLAG}: 1
\end{block}
\end{columns}

\begin{block}{Description}
Flag to switch on $\Delta-$MLFF scheme with respect to some base model.
\end{block}

\begin{block}{Remark}
Currently, the only available base model is Yukawa potential which is the default (\hyperlink{YUKAWA_A}{\texttt{YUKAWA\_A}}, \hyperlink{YUKAWA_MU}{\texttt{YUKAWA\_MU}}, \hyperlink{YUKAWA_ECONST}{\texttt{YUKAWA\_ECONST}}).
\end{block}

\end{frame}
%%%%%%%%%%%%%%%%%%%%%%%%%%%%%%%%%%%%%%%%%%%%%%%%%%%%%%%%%%%%%%%%%%%%%%%%%%%%%%%%%%%%%%%%%%%%%

%%%%%%%%%%%%%%%%%%%%%%%%%%%%%%%%%%%%%%%%%%%%%%%%%%%%%%%%%%%%%%%%%%%%%%%%%%%%%%%%%%%%%%%%%%%%%
\begin{frame}[allowframebreaks]{\texttt{{YUKAWA\_A}}} \label{YUKAWA_A}
\vspace*{-12pt}
\begin{columns}
\column{0.4\linewidth}
\begin{block}{Type}
Double array
\end{block}

\begin{block}{Default}
0.0
\end{block}

\column{0.4\linewidth}
\begin{block}{Unit}
Ha-Bohr
\end{block}

\begin{block}{Example}
\texttt{YUKAWA\_A}: 4.0 7.0 11.0
\end{block}
\end{columns}

\begin{block}{Description}
Amplitudes $A_{st}$ of the Yukawa base model $E_Y = \sum_{i<j} A_{s_i s_j}\, e^{-\mu r_{ij}}/r_{ij} + c\,N$ (sum over atom pairs and periodic images with $r_{ij} \le 15$ Bohr; $s_i$ is the type of atom $i$, $N$ the number of atoms). One value is used for every pair of atom types. Otherwise give $n(n+1)/2$ values for $n$ atom types, one per pair, in the order $A_{11}\, A_{12} \ldots A_{1n}\, A_{22} \ldots A_{2n} \ldots A_{nn}$, with the atom types numbered in the order of the \texttt{.ion} file (e.g. B, N, C: B-B B-N B-C N-N N-C C-C). Any other number of values is an error.
\end{block}

\end{frame}

\begin{frame}[allowframebreaks]{\texttt{{YUKAWA\_A}} (continued)}
\begin{block}{Remark}
In $\Delta$-MLFF with \texttt{MLFF\_FLAG}: 1, the amplitudes, \texttt{YUKAWA\_MU} and \texttt{YUKAWA\_ECONST} are fitted to the energy, forces and stress of the 5 DFT steps after \hyperlink{MLFF_SWITCH_STEP}{\texttt{MLFF\_SWITCH\_STEP}} (the input values are not used); the amplitudes and $c$ by linear least squares for each $\mu$, $\mu$ by a one-dimensional search. A pair of atom types that does not occur within 15 Bohr gets $A = 0$. The fitted values are printed in the \texttt{.out} file and written to the \texttt{.restart} file, from which a restart takes them. Only for orthogonal cells with periodic boundary conditions in all three directions.
\end{block}

\end{frame}
%%%%%%%%%%%%%%%%%%%%%%%%%%%%%%%%%%%%%%%%%%%%%%%%%%%%%%%%%%%%%%%%%%%%%%%%%%%%%%%%%%%%%%%%%%%%%

%%%%%%%%%%%%%%%%%%%%%%%%%%%%%%%%%%%%%%%%%%%%%%%%%%%%%%%%%%%%%%%%%%%%%%%%%%%%%%%%%%%%%%%%%%%%%
\begin{frame}[allowframebreaks]{\texttt{{YUKAWA\_MU}}} \label{YUKAWA_MU}
\vspace*{-12pt}
\begin{columns}
\column{0.4\linewidth}
\begin{block}{Type}
Double
\end{block}

\begin{block}{Default}
0.0
\end{block}

\column{0.4\linewidth}
\begin{block}{Unit}
Bohr$^{-1}$
\end{block}

\begin{block}{Example}
\texttt{YUKAWA\_MU}: 0.8
\end{block}
\end{columns}

\begin{block}{Description}
Inverse screening length $\mu$ of the Yukawa base model (see \hyperlink{YUKAWA_A}{\texttt{YUKAWA\_A}}), common to all pairs of atom types.
\end{block}

\end{frame}
%%%%%%%%%%%%%%%%%%%%%%%%%%%%%%%%%%%%%%%%%%%%%%%%%%%%%%%%%%%%%%%%%%%%%%%%%%%%%%%%%%%%%%%%%%%%%

%%%%%%%%%%%%%%%%%%%%%%%%%%%%%%%%%%%%%%%%%%%%%%%%%%%%%%%%%%%%%%%%%%%%%%%%%%%%%%%%%%%%%%%%%%%%%
\begin{frame}[allowframebreaks]{\texttt{{YUKAWA\_ECONST}}} \label{YUKAWA_ECONST}
\vspace*{-12pt}
\begin{columns}
\column{0.4\linewidth}
\begin{block}{Type}
Double
\end{block}

\begin{block}{Default}
0.0
\end{block}

\column{0.4\linewidth}
\begin{block}{Unit}
Ha/atom
\end{block}

\begin{block}{Example}
\texttt{YUKAWA\_ECONST}: -5.0
\end{block}
\end{columns}

\begin{block}{Description}
Constant energy per atom $c$ of the Yukawa base model (see \hyperlink{YUKAWA_A}{\texttt{YUKAWA\_A}}).
\end{block}

\end{frame}
%%%%%%%%%%%%%%%%%%%%%%%%%%%%%%%%%%%%%%%%%%%%%%%%%%%%%%%%%%%%%%%%%%%%%%%%%%%%%%%%%%%%%%%%%%%%%

%%%%%%%%%%%%%%%%%%%%%%%%%%%%%%%%%%%%%%%%%%%%%%%%%%%%%%%%%%%%%%%%%%%%%%%%%%%%%%%%%%%%%%%%%%%%%
\begin{frame}[allowframebreaks]{\texttt{{PARAM\_THERMO\_INT}}} \label{PARAM_THERMO_INT}
\vspace*{-12pt}
\begin{columns}
\column{0.4\linewidth}
\begin{block}{Type}
Double
\end{block}

\begin{block}{Default}
Not used
\end{block}

\column{0.4\linewidth}
\begin{block}{Unit}
No unit
\end{block}

\begin{block}{Example}
\texttt{PARAM\_THERMO\_INT}: 0.5
\end{block}
\end{columns}

\begin{block}{Description}
Coupling parameter $\lambda \in [0, 1]$ for thermodynamic integration from the Yukawa base model to DFT. The ions move with the forces $\lambda F_{\text{DFT}} + (1-\lambda) F_Y$, where $F_{\text{DFT}}$ are the DFT forces (or the MLFF forces when MLFF is on) and $F_Y$ those of the Yukawa model (\hyperlink{YUKAWA_A}{\texttt{YUKAWA\_A}}, \hyperlink{YUKAWA_MU}{\texttt{YUKAWA\_MU}}, \hyperlink{YUKAWA_ECONST}{\texttt{YUKAWA\_ECONST}}).
\end{block}

\end{frame}

\begin{frame}[allowframebreaks]{\texttt{{PARAM\_THERMO\_INT}} (continued)}
\begin{block}{Remark}
Each step prints the DFT, Yukawa and mixed energies and the integrand $E_{\text{DFT}} - E_Y$ in the \texttt{.out} file, and \texttt{TI\_YEN}, \texttt{TI\_MEN}, \texttt{TI\_DEN}, \texttt{TI\_TEN} in the \texttt{.aimd} file. The free energy difference is the integral over $\lambda$ from 0 to 1 of the average of \texttt{TI\_DEN} at each $\lambda$.

Needs \texttt{MD\_FLAG}: 1 with a fixed-cell MD method (NVE, NVT\_NH, NVK\_G); only the forces are mixed. Only for orthogonal cells with periodic boundary conditions in all three directions. In $\Delta$-MLFF with \texttt{MLFF\_FLAG}: 1, the forces are mixed only after the Yukawa fit (see \hyperlink{YUKAWA_A}{\texttt{YUKAWA\_A}}); until then they are the DFT forces.
\end{block}

\end{frame}
%%%%%%%%%%%%%%%%%%%%%%%%%%%%%%%%%%%%%%%%%%%%%%%%%%%%%%%%%%%%%%%%%%%%%%%%%%%%%%%%%%%%%%%%%%%%%
